% =============================================================
% Capitulo 4 - Sistemas embebidos para adquisicion y control
% =============================================================
\chapter{Sistemas embebidos para adquisici\'on y control}
\label{cap:sistemas_embebidos}

\section{Conversi\'on anal\'ogico-digital de se\~nales pulsadas}
\label{sec:adc_pulsadas}

Para digitalizar correctamente el valor retenido por el circuito
detector de picos descrito en el Cap\'itulo~\ref{cap:fotodeteccion}, el
tiempo de conversi\'on del ADC empleado debe ser lo suficientemente
corto en relaci\'on con el tiempo de descarga del capacitor de
retenci\'on, de modo de evitar que la lectura se vea afectada por dicha
descarga \cite{informeEneFeb2026}. Este criterio fue el que gui\'o la
selecci\'on del microcontrolador utilizado en el proyecto (ver
Secci\'on~\ref{sec:rp2350}).

% PLACEHOLDER: completar con el detalle de la configuracion del ADC
% integrado (resolucion, tiempo de muestreo, canales utilizados) una
% vez finalizada la etapa de implementacion del firmware.

\section{Microcontroladores de la familia RP2350}
\label{sec:rp2350}

Se seleccion\'o un RP2350 (presente en la placa de desarrollo Raspberry
Pi Pico~2) como microcontrolador del sistema, debido a que el tiempo de
conversi\'on de su ADC integrado resulta lo suficientemente corto como
para evitar la descarga significativa del circuito detector de picos
antes de completar la digitalizaci\'on, y a que el dispositivo cuenta
con documentaci\'on t\'ecnica completa \cite{informeEneFeb2026}. Para el
prototipo de la placa de control se utiliz\'o el m\'odulo comercial
Raspberry Pi Pico como referencia de dise\~no \cite{kicad_hsrl}.

% PLACEHOLDER: ampliar con arquitectura interna del RP2350 (nucleos,
% perifericos utilizados, velocidad de reloj) en base a su hoja de
% datos, a incorporar en el Anexo I.

\section{Comunicaci\'on serie RS232/UART}
\label{sec:rs232_uart}

La comunicaci\'on entre la placa de control y el l\'aser seeder se
realiza mediante el protocolo RS232 especificado por el fabricante del
equipo \cite{informeMarzo2026}. Dado que los niveles de tensi\'on RS232
no son compatibles con los niveles l\'ogicos del microcontrolador, es
necesario un circuito integrado transceptor que adapte ambos est\'andares
\cite{sat2026}. La selecci\'on de dicho componente y su integraci\'on en
la placa de control se describen en la
Secci\'on~\ref{sec:comunicacion_seeder}.
